\section{Параллелизм в обучении нейронных сетей}

Методы параллелизма, рассматриваемые в данной работе, были разработаны
для обучения больших языковых моделей (Large Language Models, LLM),
не помещающихся в память одного GPU.
Ниже описаны три основные стратегии~--- Tensor Parallelism, Pipeline
Parallelism и Fully Sharded Data Parallelism~--- и объясняется,
почему именно они рассматриваются как кандидаты для адаптации к
задаче верификации.

\subsection{Масштаб современных моделей и архитектура Transformer}

Основой современных LLM является архитектура Transformer\cite{Vaswani2017Attention},
построенная на механизме самовнимания.
Каждый слой состоит из блока Multi-Head Attention и Feed-Forward Network.

В блоке \textbf{Multi-Head Attention} вычисляются матрицы запросов,
ключей и значений $Q = XW_Q$, $K = XW_K$, $V = XW_V$,
где $W_Q, W_K, W_V \in \mathbb{R}^{d \times d}$, затем находится
взвешенная сумма:
$$\mathrm{Attention}(Q, K, V) = \mathrm{softmax}\!\left(\frac{QK^\top}{\sqrt{d}}\right)V.$$

\textbf{Feed-Forward Network}~--- двухслойная полносвязная сеть с расширением:
$$\mathrm{FFN}(x) = W_2 \cdot \mathrm{GELU}(W_1 x + b_1) + b_2,$$
где $W_1 \in \mathbb{R}^{4d \times d}$ и $W_2 \in \mathbb{R}^{d \times 4d}$.

Каждый блок содержит несколько матричных умножений с матрицами весов
размера $d \times d$ или $d \times 4d$.
При $d = 12288$ (GPT-3, 175 млрд параметров\cite{GPT3Brown2020}) одна
матрица $W_1$ занимает $12288 \times 49152 \times 4 \approx 2.3$~ГБ
в формате float32.
Модель из 96 таких слоёв не помещается ни на одном существующем
ускорителе~--- отсюда необходимость в стратегиях параллелизма.

\subsection{Tensor Parallelism (TP)}

Tensor Parallelism, предложенный в Megatron-LM\cite{MegatronLM2020},
разделяет матричные умножения \textit{внутри слоя} по нескольким GPU.
Матрица весов разрезается так, чтобы каждый GPU выполнял свою часть
умножения, а результаты синхронизировались коллективными операциями.

\subsubsection{Column Parallel Linear}

Матрица весов $W \in \mathbb{R}^{k \times n}$ разрезается по столбцам
(выходному измерению) на $P$ частей:
$$W = [W^{(0)} \mid W^{(1)} \mid \ldots \mid W^{(P-1)}], \quad
  W^{(r)} \in \mathbb{R}^{k \times (n/P)}.$$
Каждый GPU $r$ хранит только свой шард $W^{(r)}$ и вычисляет
$\mathbf{y}^{(r)} = W^{(r)\top} \mathbf{x}$, где входной вектор
$\mathbf{x}$ реплицирован на всех GPU.
После умножения выход разбит по выходному измерению.

\subsubsection{Row Parallel Linear}

Матрица весов $W \in \mathbb{R}^{n \times m}$ разбивается по строкам:
$$W = \begin{bmatrix} W^{(0)} \\ W^{(1)} \\ \vdots \\ W^{(P-1)} \end{bmatrix},
  \quad W^{(r)} \in \mathbb{R}^{(n/P) \times m}.$$
Каждый GPU $r$ умножает свою часть разбитого входа $\mathbf{x}^{(r)}$
на свой шард: $\mathbf{z}^{(r)} = W^{(r)} \mathbf{x}^{(r)}$,
затем полный выход суммируется через \texttt{AllReduce} следующим образом:
$$\mathbf{y} = \sum_{r=0}^{P-1} \mathbf{z}^{(r)}.$$

\subsubsection{Коллективные операции}

Пара Column--Row Parallel образует базовый блок TP.
На forward нужен только один \texttt{AllReduce} после Row Parallel слоя,
на backward нужен \texttt{AllReduce} для градиентов перед Column
Parallel слоем.
При быстром соединении между GPU (NVLink, NVSwitch) синхронизация
занимает сотые доли от времени вычислений\cite{MegatronLM2020}.
Без прямых соединений между GPU, например через PCIe-шину,
операция \texttt{AllReduce} становится узким местом и прирост от разбиения
сходит на нет, отсюда практическое ограничение: TP разумно
применять в рамках одной ноды, обычно не более 8~GPU.

\subsection{Pipeline Parallelism (PP)}

Pipeline Parallelism\cite{GPipePaper} разбивает модель иначе~--- не
внутри слоёв, а между ними. Каждый GPU получает свою стадию (stage):
обычно это несколько (или один) последовательных слоёв. Данные передаются через стадии как по
конвейеру, отсюда метод получил свое название. 
Можно обратить внимание, что в такой схеме неизбежно возникает проблема <<пузыря>> (pipeline bubble): пока
текущий GPU обрабатывает батч, предыдущий GPU простаивает,
на backward-pass ситуация аналогичная.

У этой проблемы есть стандартное решение~--- можно разбить входной батч на $M$ микробатчей
и запускать их один за другим без ожидания. Доля простоя в таком случае составляет:
$$r_{\text{bubble}} = \frac{P - 1}{M + P - 1},$$
при $M \gg P$ это время простоя стремится к нулю, хотя за это приходится
платить памятью для хранения промежуточных активаций всех $M$ микробатчей одновременно.
Между стадиями передаются только активации
(point-to-point send/recv), поэтому PP терпимее к медленным соединениям между GPU, чем TP.

\subsection{Fully Sharded Data Parallelism (FSDP)}

Классический Data Parallelism реплицирует всю модель на каждом GPU,
разделяя только данные и усредняя градиенты через \texttt{AllReduce}
после каждого шага. Недостаток~--- каждый GPU хранит полную копию
модели, градиентов и состояний оптимизатора, что потребляет много избыточной (redundant) памяти.

Fully Sharded Data Parallelism (FSDP), основанный на идеях статьи
ZeRO\cite{ZeROPaper}, разбивает все три компонента на части: каждый GPU хранит
лишь $1/P$ каждого тензора.
Перед вычислением конкретного слоя \texttt{AllGather} собирает
полные веса, после backward FSDP агрегирует градиенты через \texttt{ReduceScatter}
и применяет их независимо на каждом GPU к своим шардам.
Метод особенно эффективен для обучения, где состояния оптимизатора
(Adam хранит первый и второй моменты) занимают тройной объём памяти
по сравнению с весами.

\subsubsection{Стадии ZeRO и отображение на FSDP}

Статья ZeRO\cite{ZeROPaper} формализует три уровня шардирования
тренировочного состояния; каждый следующий уровень добавляет
ещё один компонент.

\textbf{ZeRO-1} разбивает только \textit{состояния оптимизатора}.
При использовании Adam каждый GPU хранит $1/P$ первых и вторых
моментов и собирает их через \texttt{AllReduce} при необходимости обновления параметров;
сами параметры и градиенты по-прежнему реплицированы на всех GPU.
Пиковая память снижается с $16\Psi$ до $(4 + 12/P)\Psi$ байт
(где $\Psi$~--- число параметров, смешанная точность).

\textbf{ZeRO-2} дополнительно шардирует \textit{градиенты}.
Вместо \texttt{AllReduce} используется \texttt{ReduceScatter}:
каждый GPU накапливает только $1/P$ усреднённых градиентов.
Параметры по-прежнему реплицированы; память $\approx (4 + 4/P)\Psi$.

\textbf{ZeRO-3} шардирует также \textit{параметры модели}~---
каждый GPU хранит ровно $1/P$ каждого весового тензора.
Перед вычислением слоя \texttt{AllGather} временно восстанавливает
полный тензор; после вычисления шарды освобождаются.
Итоговая память~--- $\approx (16/P)\Psi$, то есть экономия
пропорциональна числу GPU. Именно ZeRO-3~--- под именем
\textit{Fully Sharded Data Parallelism (FSDP)}~--- реализован
в PyTorch начиная с версии~1.12.

Для задачи верификации состояния оптимизатора и накопленные
градиенты не нужны~--- вычисление bounds это по сути инференс.
ZeRO-1 и ZeRO-2 не дают никакой экономии памяти в этом случае:
шардируемых компонентов нет. Единственная применимая стадия~---
ZeRO-3/FSDP: параметры шардированы, перед каждым слоем
\texttt{AllGather} собирает полные веса, после обработки шарды
освобождаются~--- и всё это без каких-либо градиентных операций.

Именно FSDP стал основным алгоритмом адаптации в данной работе:
при верификации состояний оптимизатора нет, а значит \texttt{AllGather}/free
шардов можно встроить в процедуру bound propagation без изменения
логики backward-pass~--- подробнее в главе~\ref{sec:parallel-verifiers}.
